%!TEX root = ../../main.tex
\section{적정 점수, 보정, 짝 비교}
\label{sec:rel-scores}

제 \ref{sec:intro-approach} 절에서 예측이 틀린 정도를 Brier 점수로 재고, 두 예측을 같은 교전끼리 견주며, 경기를 다시 뽑아 그 차이가 놓이는 범위를 구한다고 하였다. 이 절은 그 도구들을 정의한다. 확률 예측을 채점하는 점수, 점수를 성분으로 나누는 방법, 순위를 보는 보조 지표, 확률을 다시 맞추는 보정, 두 예측의 점수 차이를 견주는 방법의 순서이다.

먼저 확률 예측을 채점하는 점수이다. 본 논문의 예측은 추정 승률이 오를 확률이고 결과는 올랐는지 여부이므로, 확률과 이진 결과 사이의 거리를 재는 점수가 필요하다.

\begin{definition}[Brier 점수와 로그 손실]
확률 예측 $p_i$와 이진 결과 $y_i \in \{0,1\}$, 가중치 $w_i$에 대해 가중 Brier 점수는
\begin{equation}
  \Brier = \frac{\sum_i w_i (p_i - y_i)^2}{\sum_i w_i}
  \label{eq:brier}
\end{equation}
이고, 로그 손실은
\begin{equation}
  -\frac{\sum_i w_i \bigl[ y_i \log p_i + (1 - y_i) \log (1 - p_i) \bigr]}{\sum_i w_i}
  \label{eq:logloss}
\end{equation}
이다. Brier 점수는 1950년 Brier가 기상 예보의 검증을 위해 제안하였다 \cite{brier1950verification}. 두 점수 모두 엄격한 적정 채점 규칙(strictly proper scoring rule), 즉 예보자가 자신의 참 확률을 보고할 때에만 기대 점수가 최소가 되는 규칙이며, 그 일반 이론은 Gneiting과 Raftery \cite{GneitingRaftery2007}가 정립하였다.
\end{definition}

두 점수 모두 낮을수록 덜 틀린 것이다. 모든 사례에 0.5라고 답하면 Brier 점수는 0.25이다. 적정 채점 규칙이라는 성질은 확률을 실제보다 크거나 작게 말해서는 점수를 좋게 만들 수 없다는 뜻이며, 이 성질을 갖는 점수를 이하 적정 점수라 부른다. 본 논문은 Brier 점수를 주 점수로 쓴다.

다음은 Brier 점수의 차이가 어디서 나오는지를 나누어 보는 방법이다.

\begin{definition}[Brier 점수의 분해와 CORP]
Brier 점수는 세 성분으로 나뉜다. 말한 확률이 실제 비율과 어긋난 정도인 오보정(miscalibration) 성분 $\MCB$, 결과가 다른 사례에 다른 확률을 주어 구별한 정도인 판별(discrimination) 성분 $\DSC$, 결과 자체가 갖는 변동인 불확실성(uncertainty) 성분 $\UNC$이다. 이 분해는 Murphy \cite{murphy1973vector}가 제안하였고, 구간 나누기에 의존하지 않는 안정적 계산법으로 Dimitriadis, Gneiting, Jordan \cite{DimitriadisGneitingJordan2021}이 CORP(consistent, optimally binned, reproducible, PAV-based) 분해를 제안하였다. 평가 표본에서 같은 가중치 $w_i$를 쓴 가중 단조 회귀(isotonic regression)로 재보정한 예측 $p^{*}$의 가중 Brier 점수를 $\Brier_{\mathrm{iso}}$라 하면
\begin{equation}
\begin{aligned}
  \MCB &= \Brier - \Brier_{\mathrm{iso}}, \\
  \DSC &= \UNC - \Brier_{\mathrm{iso}}, \\
  \Brier &= \MCB - \DSC + \UNC.
\end{aligned}
\label{eq:corp}
\end{equation}
$\UNC$는 결과의 가중 양성 비율 $\bar y_w$만으로 정해진다.
\begin{equation}
  \bar y_w = \frac{\sum_i w_i y_i}{\sum_i w_i}, \qquad \UNC = \bar y_w (1 - \bar y_w).
  \label{eq:unc}
\end{equation}
\end{definition}

$\MCB$는 작을수록, $\DSC$는 클수록 좋다. $\UNC$는 예측과 무관하게 결과의 가중 양성 비율만으로 정해지며, 모든 사례에 그 비율을 답하는 예측의 가중 Brier 점수와 같다. 본 논문에서 이 분해는 평가 자료에서 두 예측의 점수 차이가 어느 성분에서 나오는지를 보는 진단으로만 쓴다. 단조 회귀는 분해를 계산하는 장치이며, 보고하는 모델에는 적용하지 않는다.

Brier 점수는 확률의 값 자체를 채점한다. 사례의 순위만 보는 보조 지표로는 AUC를 쓴다.

\begin{definition}[AUC]
점수 함수 $f$에 대해 AUC는 무작위로 뽑은 양성 표본의 점수가 음성 표본의 점수보다 클 확률이며, Mann--Whitney 통계량과 동치임을 Hanley와 McNeil \cite{hanley1982meaning}이 보였다. AUC $0.64$는 무작위 양성 교전의 점수가 무작위 음성 교전의 점수보다 클 확률이 0.64라는 뜻이다.
\end{definition}

본 논문은 AUC를 순위가 얼마나 맞는지를 보는 보조 지표로 보고하고, 주된 판단은 Brier 점수의 차이로 한다.

\begin{definition}[사후 보정]
모형이 낸 확률을 단조 함수에 넣어 실제 비율에 맞도록 다시 맞추는 절차를 사후 보정이라 한다. Platt 척도 \cite{platt2000probabilities}, 온도 척도 \cite{guo2017calibration}, 구간 기반 보정 \cite{naeini2015obtaining}이 표준적인 도구이다.
\end{definition}

본 논문이 사후 보정을 어떻게 적용하였는지는 승률 모형에 대해서는 제 \ref{sec:value-evaluator} 절에, 교전 전 상태의 모델과 기준선에 대해서는 제 \ref{sec:pred-selection} 절에 있다.

마지막은 두 예측의 점수 차이를 견주는 방법이다.

\begin{definition}[짝 비교와 군집 부트스트랩]
같은 사례에 대한 두 예측 $A$, $B$의 점수 차이는
\begin{equation}
  \dBrier(A - B) = \Brier(A) - \Brier(B)
  \label{eq:dbrier}
\end{equation}
이다. 음수이면 $A$가 덜 틀린 것이다. 같은 사례에서 얻은 두 확률 예측은 이렇게 짝지어 견줄 수 있다. 다만 한 경기에서 나온 교전들은 서로 닮아 있어 독립된 표본으로 볼 수 없다. 그래서 본 논문은 교전이 아니라 경기를 단위로 복원 추출하고, 뽑힌 경기의 교전을 모두 가져와 $\dBrier$를 다시 계산하는 일을 반복하는 짝 경기 군집 부트스트랩 \cite{efron1993bootstrap,CameronMiller2015}으로 $\dBrier$의 95\% 백분위 구간을 구한다.
\end{definition}

학습 알고리즘의 선택에 관해서도 적어 둔다. 표 형태의 자료에서는 트리 앙상블과 심층 신경망 가운데 어느 쪽이 나은지가 벤치마크마다 달라진다 \cite{grinsztajn2022tree,gorishniy2021revisiting,shwartzziv2022tabular,ke2017lightgbm}. 본 논문은 교전 전 상태의 모델을 로지스틱 회귀 하나로 미리 정해 두었고, 입력의 종류를 견주는 실험에서만 트리 앙상블인 LightGBM을 썼다(제 \ref{sec:pred-infogroups} 절).
